% =========================================================
% TABLA 6: OBJETIVO ESPECÍFICO 6
% =========================================================
\begin{table}[htbp]
	\centering
	\caption{Cumplimiento del Objetivo Específico 6}
	\label{tab:cumplimiento-objetivo-6}
	\begin{tabular}{|>{\raggedright\arraybackslash}p{0.38\textwidth}|>{\raggedright\arraybackslash}p{0.56\textwidth}|}
		\hline
		\rowcolor[gray]{0.85}
		\multicolumn{2}{|c|}{\textbf{\textit{Objetivo específico 6}}}                                                                                                                                                                                                                                                                                                           \\
		\hline
		\multicolumn{2}{|>{\centering\arraybackslash}p{\dimexpr0.94\textwidth+2\tabcolsep+\arrayrulewidth\relax}|}{\textit{Validar prototipo funcional del servicio de directorio implementado para el GRID}}                                                                                                                                                                                \\
		\hline
		\textbf{\textit{Detalle de cumplimiento}}                                                                                                                                                                                                                                                                                                          & \textbf{\textit{Observaciones}} \\
		\hline
		Ver capítulo \textcolor[HTML]{1B6B93}{\textbf{\ref{cap:pruebas}}} \textbf{Pruebas}. Se ejecutaron nueve pruebas que abordaron la integración, autenticación centralizada, autorización diferenciada, gestión del ciclo de vida de cuentas, políticas de seguridad, trazabilidad, replicación, continuidad ante fallas y recuperación del servicio. &
		Las pruebas ejecutadas fueron (sección \textcolor[HTML]{1B6B93}{\ref{sec:ejecucion-pruebas}}):\par\vspace{0.4em}
		- \textbf{PT-01}: Integración y autenticación centralizada (sección \textcolor[HTML]{1B6B93}{\ref{subsec:pt01}}).\par
		- \textbf{PT-02}: Autorización diferenciada mediante grupos (sección \textcolor[HTML]{1B6B93}{\ref{subsec:pt02}}).\par
		- \textbf{PT-03}: Gestión del ciclo de vida de las cuentas (sección \textcolor[HTML]{1B6B93}{\ref{subsec:pt03}}).\par
		- \textbf{PT-04}: Aplicación de políticas de seguridad (sección \textcolor[HTML]{1B6B93}{\ref{subsec:pt04}}).\par
		- \textbf{PT-05}: Registro y trazabilidad (sección \textcolor[HTML]{1B6B93}{\ref{subsec:pt05}}).\par
		- \textbf{PT-06}: Revisión de permisos (sección \textcolor[HTML]{1B6B93}{\ref{subsec:pt06}}).\par
		- \textbf{PT-07}: Replicación y sincronización (sección \textcolor[HTML]{1B6B93}{\ref{subsec:pt07}}).\par
		- \textbf{PT-08}: Continuidad ante falla (sección \textcolor[HTML]{1B6B93}{\ref{subsec:pt08}}).\par
		- \textbf{PT-09}: Recuperación y resincronización (sección \textcolor[HTML]{1B6B93}{\ref{subsec:pt09}}).\par\vspace{0.4em}
		Las limitaciones encontradas (autenticación multifactor y recuperación de contraseñas) quedaron documentadas como trabajo futuro, y la evaluación consolidada del cumplimiento de requisitos se presenta en la \autoref{tab:evaluacion-requisitos}.                                                                                                                                  \\
		\hline
	\end{tabular}
\end{table}
